
\section[\trSymplecticDiffeomorphism{}与\trHamiltonianVectorField]{$(\R^{2n},\omega_0)$ 中的\trSymplecticDiffeomorphism{}与\trHamiltonianVectorField}
\label{sec:symplectic-diffeomorphisms}

现在，我们转向\trSymplectic{}空间 $(\R^{2n},\omega_0)$ 中的\trNonlinear{}映射。%
$\R^{2n}$ 中的一个\textbf{\trDiffeomorphism}（diffeomorphism）$\varphi$ 称为\trSymplectic{}的，若
\begin{equation}\label{eq:symplectic-diffeomorphism}
\varphi^{*}\omega_0=\omega_0,
\end{equation}
其中按定义，任意\trTwoForm{} $\omega$ 的\trPullback{}由
\[
(\varphi^{*}\omega)_x(a,b)=\omega_{\varphi(x)}\bigl(\varphi'(x)a,\ \varphi'(x)b\bigr)
\]
给出，这里 $x\in\R^{2n}$，而 $a$, $b\in T_{x}\R^{2n}=\R^{2n}$。%
此处 $\varphi'(x)$ 表示 $\varphi$ 在点 $x$ 处的\textbf{\trDerivative}（derivative），由\trJacobi{}矩阵表出。%
鉴于 $\omega_0$ 的定义，$(\R^{2n},\omega_0)$ 中的%
\textbf{\trSymplecticDiffeomorphism}（symplectic diffeomorphism）便由 $\varphi$ 的一阶\trDerivative{}所满足的恒等式
\begin{equation}\label{eq:first-derivative-identity}
\varphi'(x)^{T}J\varphi'(x)=J,\qquad x\in\R^{2n}
\end{equation}
所刻画。于是 $\varphi'(x)$ 是一个\trSymplecticMatrix，特别地，
\begin{equation}\label{eq:det-first-derivative}
\det\varphi'(x)=1,
\end{equation}
从而\trSymplecticDiffeomorphism{}是\trVolumePreserving{}的。然而，当 $n>1$ 时，%
\trSymplecticDiffeomorphism{}的类远较\trVolumePreserving{}\trDiffeomorphism{}受限。%
这一点在下文将会明朗：届时我们循着\trGromov{}的思路考察这样一个问题——%
用\trSymplecticDiffeomorphism{}究竟能做些什么？

$(\R^{2n},\omega_0)$ 中的\trSymplecticDiffeomorphism{} $\varphi$ 不仅保持 $\omega_0$ 与相应的\trVolumeForm{} $\Omega$，%
而且还保持\textbf{\trClosed\trCurve}（closed curve）的\textbf{\trAction}（action），这一点我们接下来就会看到。%
形式 $\omega_0$ 是\textbf{\trExact{}形式}（exact form），因为
\begin{equation}\label{eq:omega0-exact}
\omega_0=\sum_{j=1}^{n}dy_j\wedge dx_j=d\lambda,
\end{equation}
其中\textbf{\trOneForm}（1-form）$\lambda$ 由下式定义：
\[
\lambda=\sum_{j=1}^{n}y_j\,dx_j.
\]
因此，只要 $\varphi$ 是\trSymplectic{}的，$\lambda-\varphi^{*}\lambda$ 便是\textbf{\trClosed{}形式}（closed form）。%
事实上，$d(\lambda-\varphi^{*}\lambda)=d\lambda-d\bigl(\varphi^{*}\lambda\bigr)=d\lambda-\varphi^{*}d\lambda=\omega_0-\varphi^{*}\omega_0=0$。%
利用\trPoincare{}引理可以找到一个函数 $F:\R^{2n}\to\R$ 满足
\begin{equation}\label{eq:lambda-dF}
\lambda-\varphi^{*}\lambda=dF.
\end{equation}
若 $\gamma$ 是一条\trOriented{}的\textbf{简单\trClosed\trCurve}（simply closed curve），则可以积分，并借助 \eqref{eq:lambda-dF} 得到
\[
\int_{\gamma}\lambda=\int_{\gamma}\varphi^{*}\lambda=\int_{\varphi(\gamma)}\lambda
\]
这是因为\trExact{}形式沿\trClosed\trCurve{}的积分必为零。若把\trClosed\trCurve{} $\gamma$ 的\trAction{} $A(\gamma)$ 定义为
\begin{equation}\label{eq:action-definition}
A(\gamma)=\int_{\gamma}\lambda\in\R,
\end{equation}
便看到
\begin{equation}\label{eq:action-invariance}
A\bigl(\varphi(\gamma)\bigr)=A(\gamma)
\end{equation}
（只要 $\varphi$ 是\trSymplectic{}的）；于是如所断言，$\varphi$ 保持\trAction{}不变。%
反过来，当然，若 $\R^{2n}$ 中的一个\trDiffeomorphism{} $\varphi$ 对 $\R^{2n}$ 中的一切\trClosed\trCurve{}%
都满足 \eqref{eq:action-invariance}，我们便得出 $\varphi$ 是\trSymplectic{}的结论。%
若用 $x(t)$ 将 $\gamma$ \trParametrization（$0\leq t\leq1$，$x(0)=x(1)$），则\trAction{}便成为
\begin{equation}\label{eq:action-integral}
A(\gamma)=\frac{1}{2}\int_{0}^{1}\langle-J\dot{x},x\rangle\,dt.
\end{equation}

\trSymplecticDiffeomorphism{}的例子由所谓的\trHamiltonianVectorField{}产生，我们现在对此作一回顾。%
对于\trSymplecticForm{} $\omega_0$ 与一个\trSmooth{}函数
\[
H:\R^{2n}\to\R,
\]
我们可以在 $\R^{2n}$ 上配以一个\textbf{\trVectorField}（vector field）$X_H$，即要求
\begin{equation}\label{eq:hamiltonian-vector-field}
\omega_0\bigl(X_H(x),a\bigr)=-dH(x)a
\end{equation}
对一切 $a\in\R^{2n}$ 与 $x\in\R^{2n}$ 成立。由于 $\omega_0$ \trNondegenerate，向量 $X_H(x)$ 便唯一确定。%
条件 \eqref{eq:hamiltonian-vector-field} 等价于 $\langle JX_H(x),a\rangle=-\langle\nabla H(x),a\rangle$，%
其中 $H$ 的\textbf{\trGradient}（gradient）一如通常地相对于\trEuclidean\trInnerProduct{}来定义。%
因此 $JX_H(x)=-\nabla H(x)$，并且由 $J^{2}=-1$ 便得到表示
\begin{equation}\label{eq:xh-representation}
X_H(x)=J\nabla H(x),\qquad x\in\R^{2n}.
\end{equation}

显然，\trHamiltonianVectorField{}是非常特殊的。它们尤其与\trGradient{}型 $X=\nabla H(x)$ 的向量场截然不同，因为 $J$ 是\trAntisymmetric{}的。

下文中，我们用 $\varphi^{t}$ 表示向量场 $X$ 的\textbf{\trFlow}（flow）。它由
\begin{align*}
\frac{d}{dt}\varphi^{t}(x) &= X\bigl(\varphi^{t}(x)\bigr)\\
\varphi^{0}(x) &= x,\qquad x\in\R^{2n}.
\end{align*}
所定义。\trCurve{} $x(t)=\varphi^{t}(x)$ 对于初始条件 $x\in\R^{2n}$ 求解\trCauchy{}\textbf{\trInitialValueProblem}（initial value problem）。%
现在设 $X=X_H$ 是由 $\omega_0$ 与 $H$ 确定的\trHamiltonianVectorField。那么每一个流映射 $\varphi^{t}$ 都保持形式 $\omega_0$：
\begin{equation}\label{eq:flow-preserves-form}
\bigl(\varphi^{t}\bigr)^{*}\omega_0=\omega_0,
\end{equation}
从而是一个\trSymplecticMapping。这容易验证，下一节将在更抽象的框架中给出证明。

为了下文之用，回顾一下 $\R^{m}$ 上向量场 $X$ 的变换公式是有益的。设 $x(t)$ 是\textbf{\trDifferentialEquation}（differential equation）
\[
\dot{x}=X(x),\qquad x\in\R^{m}.
\]
的解。若 $u:\R^{m}\to\R^{m}$ 是一个\trDiffeomorphism，我们便可借
\[
x(t)=u\bigl(y(t)\bigr)
\]
来定义\trCurve{} $y(t)$。对 $t$ 求导我们便得出：$y(t)$ 求解方程
\[
\dot{y}=Y(y),\qquad y\in\R^{m}
\]
其中 $Y$ 是由
\[
Y(y)=du(y)^{-1}\cdot X\circ u(y).
\]
所定义的变换后的向量场。下文中我们将使用记号
\begin{equation}\label{eq:pullback-vector-field}
u^{*}X:=(du)^{-1}\cdot X\circ u.
\end{equation}
我们已经证明：$X$ 的\trFlow{} $\varphi^{t}$ 与 $u^{*}X$ 的\trFlow{} $\psi^{t}$ 这两个\trFlow{}%
由\trDiffeomorphism{} $u$ 所\trConjugate，即
\[
\varphi^{t}\circ u=u\circ\psi^{t}.
\]
如果让 $\R^{2n}$ 中的一个\trHamiltonianVectorField{} $X_H$ 经受任意的变换 $u$，%
其特殊形式便会被破坏。然而，\trSymplectic{}变换却保持\trHamiltonianVectorField{}这一类。%
事实上，若 $u^{*}\omega_0=\omega_0$，则
\begin{equation}\label{eq:sympl-change-hamiltonian}
u^{*}X_H=X_K\quad\text{且}\quad K=H\circ u.
\end{equation}
这容易验证：把函数 $K$ 定义为复合 $K=H\circ u$，则由\textbf{\trChainRule}（chain rule），$dK=dH\circ u\cdot du$，%
而相对于\trEuclidean\trInnerProduct{}取的\trGradient{}便成为 $\nabla K=(du)^{T}\nabla H\circ u$。%
由假设，$du$ 在每一点都是\trSymplecticMapping，从而 $(du)^{T}$ 亦然，于是 $du\cdot J\cdot(du)^{T}=J$。%
于是，由\trHamiltonianVectorField{}的定义 \eqref{eq:xh-representation}，
\begin{align*}
X_K &= J\nabla K = J(du)^{T}\nabla H\circ u\\
    &= (du)^{-1}(J\nabla H)\circ u\\
    &= u^{*}X_H,
\end{align*}
这正是所要证明的。
